% !TeX root = ../apunte_mecanica_clasica.tex
\IfFileExists{../apunte_mecanica_clasica.tex}{%
  \documentclass[../apunte_mecanica_clasica.tex]{subfiles}%
}{%
  \documentclass[apunte_mecanica_clasica.tex]{subfiles}%
}
\begin{document}

\chapter{Fuerzas centrales, ecuación de Binet y problema de Kepler}
\label{ch:fuerzas-centrales-binet-kepler}
% Figuras extraídas desde Clase MC 2-2.pdf y organizadas en figures/09_fuerzas_centrales_binet_kepler/.

\section{Fuerzas centrales}

\subsection{Definición de fuerza central}

Una fuerza $\vec F$ se dice que es una \textbf{fuerza central} si la dirección de ésta pasa por el punto donde está el centro de la fuerza.

La definición anterior establece que si un cuerpo, ubicado en la posición $P$, está sometido a la acción de un agente externo ubicado en un punto $Q$, entonces la fuerza central ejercida por el agente externo es proporcional al vector que va del punto $Q$ al punto $P$, esto es
\begin{equation}
    \vec F \propto \vec x = \vec x_P - \vec x_Q .
    \label{eq:fcbk-fuerza-central-proporcional}
\end{equation}
La construcción geométrica asociada a esta definición se muestra en la Figura~\ref{fig:fcbk-fuerza-central-geometria}.

\begin{figure}[H]
    \centering
    \IfFileExists{figures/09_fuerzas_centrales_binet_kepler/fuerza_central_geometria.png}{%
        \includegraphics[width=0.55\linewidth]{figures/09_fuerzas_centrales_binet_kepler/fuerza_central_geometria.png}%
    }{%
        \fbox{\parbox{0.55\linewidth}{\centering Figura pendiente: geometría de una fuerza central.}}%
    }
    \caption{Geometría de una fuerza central: el agente externo se ubica en $Q$ y el cuerpo en $P$.}
    \label{fig:fcbk-fuerza-central-geometria}
\end{figure}

Por comodidad, ajustamos el sistema de coordenado tal que $\vec x_Q=\vec 0$, entonces $\vec x=\vec x_P$, es decir, lo colocamos sobre $Q$. Por lo tanto,
\begin{equation}
    \vec F = f(x^i)\,\vec x,
    \label{eq:fcbk-fuerza-central-general}
\end{equation}
donde $f(x^i)$ es el factor de proporcionalidad y puede depender de las coordenadas. Se satisface:
\begin{itemize}
    \item si $f(x^i)<0$, la fuerza es \textbf{atractiva};
    \item si $f(x^i)>0$, la fuerza es \textbf{repulsiva}.
\end{itemize}

\subsection{Conservación del momentum angular}

\begin{quote}
\textbf{Proposición.} En un movimiento central el \textbf{momentum angular} es constante.
\end{quote}

\noindent
\textbf{Demostración.} Tenemos que el torque está dado por
\begin{equation}
    \begin{aligned}
        \frac{d\vec l}{dt}
        &= \vec\tau \\
        &= \vec x \times \vec F \\
        &= \vec x \times f(x^i)\vec x \\
        &= f(x^i)\,\vec x\times \vec x \\
        &= \vec 0 .
    \end{aligned}
    \label{eq:fcbk-torque-fuerza-central}
\end{equation}
Luego,
\begin{equation}
    \frac{d\vec l}{dt}=\vec 0,
    \qquad \Longrightarrow \qquad
    \vec l=\text{constante}.
    \label{eq:fcbk-momentum-angular-constante}
\end{equation}
Por lo tanto, $\vec l$ es una constante del movimiento. \hfill $\square$

Puesto que $\vec l$ es un vector, o sea posee tres componentes, la proposición anterior proporciona tres constantes del movimiento: $l_x$, $l_y$ y $l_z$. Además, el módulo de $\vec l$ también es constante.

\subsection{El movimiento central es planar}

\begin{quote}
\textbf{Proposición.} Un movimiento central es un \textbf{movimiento planar}.
\end{quote}

\noindent
\textbf{Demostración.} Puesto que $\vec l$ es constante, en magnitud y dirección, y como sabemos que el momentum angular de una partícula, que posee momentum lineal $\vec p=m\vec v$, con respecto al eje de giro que pasa por $O$, es dado por
\begin{equation}
    \vec l=\vec x\times \vec p = m\vec x\times \vec v .
    \label{eq:fcbk-def-momentum-angular}
\end{equation}
Dado que $\vec x$ y $\vec v$ definen un plano, entonces $\vec l$ es perpendicular al plano formado por $\vec x$ y $\vec v$. Podemos concluir que como $\vec l$ es constante entonces $\vec x$ y $\vec v$ están permanentemente en un plano. Por lo tanto, un movimiento central es un movimiento planar, como se ilustra en la Figura~\ref{fig:fcbk-movimiento-central-planar}. \hfill $\square$

\begin{figure}[H]
    \centering
    \IfFileExists{figures/09_fuerzas_centrales_binet_kepler/movimiento_central_planar.png}{%
        \includegraphics[width=0.55\linewidth]{figures/09_fuerzas_centrales_binet_kepler/movimiento_central_planar.png}%
    }{%
        \fbox{\parbox{0.55\linewidth}{\centering Figura pendiente: movimiento central planar.}}%
    }
    \caption{El vector $\vec l$ es perpendicular al plano generado por $\vec x$ y $\vec v$.}
    \label{fig:fcbk-movimiento-central-planar}
\end{figure}

De lo anterior, podemos ajustar el sistema coordenado de tal forma que hacemos $z=0$ para todo el movimiento. Por lo tanto,
\begin{equation}
    \vec F = f(x,y)\vec x = f(x,y)(x\hat{\imath}+y\hat{\jmath}).
    \label{eq:fcbk-fuerza-plano-xy}
\end{equation}
Esto significa que la coordenada $z$ se elimina de las ecuaciones del movimiento, como se representa en la Figura~\ref{fig:fcbk-plano-z0}.

\begin{figure}[H]
    \centering
    \IfFileExists{figures/09_fuerzas_centrales_binet_kepler/plano_movimiento_z0.png}{%
        \includegraphics[width=0.55\linewidth]{figures/09_fuerzas_centrales_binet_kepler/plano_movimiento_z0.png}%
    }{%
        \fbox{\parbox{0.55\linewidth}{\centering Figura pendiente: elección del plano $z=0$.}}%
    }
    \caption{Elección del plano de movimiento tal que $z=0$.}
    \label{fig:fcbk-plano-z0}
\end{figure}

La elección del sistema coordenado con $z=0$ implica que dos componentes del momentum angular son cero, $l_x=0$ y $l_y=0$, entonces
\begin{equation}
    \vec l = l_z\hat{k}.
    \label{eq:fcbk-momentum-angular-z}
\end{equation}
Así, sólo hay una constante del movimiento asociada al momentum angular.

\subsection{Fuerza central conservativa}

Si exigimos que la fuerza central sea conservativa entonces se debe cumplir que
\begin{equation}
    \vec\nabla\times \vec F=\vec 0.
    \label{eq:fcbk-condicion-conservativa}
\end{equation}
En el plano $z=0$, usando \eqref{eq:fcbk-fuerza-plano-xy}, se obtiene
\begin{equation}
    \begin{aligned}
    \vec\nabla\times\vec F
    &=
    \begin{vmatrix}
    \hat{\imath} & \hat{\jmath} & \hat{k} \\
    \partial_x & \partial_y & \partial_z \\
    x f(x,y) & y f(x,y) & 0
    \end{vmatrix} \\
    &=\left(y\frac{\partial f}{\partial x}-x\frac{\partial f}{\partial y}\right)\hat{k}.
    \end{aligned}
    \label{eq:fcbk-racional-fuerza-central}
\end{equation}
Por lo tanto, la condición conservativa implica
\begin{equation}
    y\frac{\partial f}{\partial x}-x\frac{\partial f}{\partial y}=0.
    \label{eq:fcbk-condicion-fxy}
\end{equation}

La expresión anterior se puede escribir en la forma
\begin{equation}
    y\frac{\partial f}{\partial x}-x\frac{\partial f}{\partial y}
    =-(-y,x,0)\cdot \vec\nabla f =0,
    \label{eq:fcbk-vector-tangente-f}
\end{equation}
donde $(-y,x,0)$ es un vector tangente a la circunferencia de radio $\sqrt{x^2+y^2}$.

Introducimos, por la geometría del problema, las coordenadas cilíndricas:
\begin{equation}
    x=r\cos\phi,
    \qquad
    y=r\sin\phi,
    \qquad
    z=z,
    \label{eq:fcbk-coordenadas-cilindricas}
\end{equation}
con $r>0$, $0\leq \phi<2\pi$ y $z\in\mathbb{R}$. Los vectores base coordenado son
\begin{equation}
    \vec e_r=\cos\phi\,\hat{\imath}+\sin\phi\,\hat{\jmath}=\hat r,
    \qquad
    \vec e_\phi=r(-\sin\phi\,\hat{\imath}+\cos\phi\,\hat{\jmath})=r\hat\phi,
    \qquad
    \vec e_z=\hat{k}.
    \label{eq:fcbk-bases-cilindricas}
\end{equation}
Así tenemos que
\begin{equation}
    (-y,x,0)=(-r\sin\phi,r\cos\phi,0)=\vec e_\phi=r\hat\phi,
    \label{eq:fcbk-vector-tangente-cilindrico}
\end{equation}
y el gradiente se escribe como
\begin{equation}
    \vec\nabla f = \vec e_r\frac{\partial f}{\partial r}
    + \vec e_\phi\frac{1}{r^2}\frac{\partial f}{\partial \phi}.
    \label{eq:fcbk-gradiente-f-cilindrico}
\end{equation}
Entonces,
\begin{equation}
    y\frac{\partial f}{\partial x}-x\frac{\partial f}{\partial y}
    =\frac{\partial f(r,\phi)}{\partial \phi}=0.
    \label{eq:fcbk-f-no-depende-phi}
\end{equation}
Por lo tanto,
\begin{equation}
    f(r,\phi)=f(r),
    \label{eq:fcbk-f-de-r}
\end{equation}
es decir, no depende del ángulo $\phi$.

Como $\vec x=r\hat r$, la fuerza central queda
\begin{equation}
    \vec F(r,\phi)=f(r)\vec x=f(r)r\hat r.
    \label{eq:fcbk-fuerza-fr-rhat-intermedia}
\end{equation}
Sin pérdida de generalidad, absorbemos la variable $r$ en $f(r)$, o sea hacemos $rf(r)\to f(r)$, aunque puedan poseer distintas dimensionalidades. Tenemos así
\begin{equation}
    \vec F(r,\phi)=f(r)\hat r.
    \label{eq:fcbk-fuerza-central-radial}
\end{equation}

\subsection{Resumen de la reducción planar}

Resumiendo lo obtenido hasta ahora: en una fuerza central
\begin{equation}
    \vec F(x^i)=f(x^i)\vec x,
    \label{eq:fcbk-fuerza-central-resumen-1}
\end{equation}
el momentum angular se conserva. Como consecuencia, el movimiento es planar, por lo que se ajusta el sistema de coordenado, con el fin de simplificar las ecuaciones del movimiento, tal que $z=0$.

Si se exige que la fuerza central sea conservativa, se concluye que
\begin{equation}
    f(x,y)=f\left(\sqrt{x^2+y^2}\right).
    \label{eq:fcbk-f-radial-cartesiano}
\end{equation}
Cambiando a coordenadas cilíndricas, la expresión de la fuerza central cambia a
\begin{equation}
    \vec F(r,\phi)=f(r)\hat r,
    \label{eq:fcbk-fuerza-central-resumen-2}
\end{equation}
y el momentum angular toma la forma
\begin{equation}
    \vec l=l_z\hat{k}=m(x\dot y-y\dot x)\hat{k}=mr^2\dot\phi\,\hat{k}.
    \label{eq:fcbk-l-mr2phidot-vector}
\end{equation}
Por lo tanto,
\begin{equation}
    l=mr^2\dot\phi
    \label{eq:fcbk-l-mr2phidot}
\end{equation}
es constante.

\newpage
\section{Ecuación de Binet y rapidez areolar}

\subsection{Ecuaciones de movimiento en coordenadas cilíndricas}

La aceleración en coordenadas cilíndricas está dada por
\begin{equation}
    \vec a=(\ddot r-r\dot\phi^2)\hat r+(2\dot r\dot\phi+r\ddot\phi)\hat\phi.
    \label{eq:fcbk-aceleracion-cilindrica}
\end{equation}
Por la segunda ley de Newton, usando \eqref{eq:fcbk-fuerza-central-resumen-2}, tenemos
\begin{equation}
    m\left[(\ddot r-r\dot\phi^2)\hat r+(2\dot r\dot\phi+r\ddot\phi)\hat\phi\right]
    =f(r)\hat r+0\hat\phi.
    \label{eq:fcbk-newton-cilindricas}
\end{equation}
Comparando componentes se concluye que
\begin{align}
    m(\ddot r-r\dot\phi^2)&=f(r),
    \label{eq:fcbk-ecuacion-radial}\\
    m(2\dot r\dot\phi+r\ddot\phi)&=0.
    \label{eq:fcbk-ecuacion-angular}
\end{align}
De la segunda ecuación se tiene
\begin{equation}
    2\dot r\dot\phi+r\ddot\phi
    =\frac{1}{r}\left(2r\dot r\dot\phi+r^2\ddot\phi\right)
    =\frac{1}{r}\frac{d}{dt}\left(r^2\dot\phi\right).
    \label{eq:fcbk-derivada-r2phidot}
\end{equation}
Entonces, de \eqref{eq:fcbk-ecuacion-angular} y \eqref{eq:fcbk-derivada-r2phidot}, se obtiene
\begin{equation}
    r^2\dot\phi=\text{constante}.
    \label{eq:fcbk-r2phidot-constante}
\end{equation}
Esto es equivalente a la conservación del momentum angular escrita en \eqref{eq:fcbk-l-mr2phidot}.

\subsection{Rapidez areolar}

Consideremos dos posiciones de la partícula en el instante $t$ y $t+dt$. El área que describe el vector posición, cuando se mueve desde $Q$ a $Q'$, es
\begin{equation}
    dA=\frac{1}{2}bdh.
    \label{eq:fcbk-area-diferencial-triangulo}
\end{equation}
Con $b=r+dr$ y $dh=r\sin(d\phi)\simeq r\,d\phi$, entonces
\begin{equation}
    dA=\frac{1}{2}(r+dr)r\,d\phi\simeq \frac{1}{2}r^2d\phi.
    \label{eq:fcbk-dA-areolar}
\end{equation}
De esta manera, la \textbf{rapidez areolar} queda dada por
\begin{equation}
    v_a=\frac{dA}{dt}=\frac{1}{2}r^2\frac{d\phi}{dt}=\frac{1}{2}r^2\dot\phi.
    \label{eq:fcbk-rapidez-areolar}
\end{equation}
La interpretación geométrica del área barrida por el radio vector se muestra en la Figura~\ref{fig:fcbk-rapidez-areolar}.

\begin{figure}[H]
    \centering
    \IfFileExists{figures/09_fuerzas_centrales_binet_kepler/rapidez_areolar.png}{%
        \includegraphics[width=0.62\linewidth]{figures/09_fuerzas_centrales_binet_kepler/rapidez_areolar.png}%
    }{%
        \fbox{\parbox{0.62\linewidth}{\centering Figura pendiente: rapidez areolar.}}%
    }
    \caption{Área diferencial barrida por el radio vector durante el intervalo $dt$.}
    \label{fig:fcbk-rapidez-areolar}
\end{figure}

Por \eqref{eq:fcbk-r2phidot-constante}, se obtiene
\begin{equation}
    r^2\dot\phi=\text{constante},
    \qquad
    2v_a=\text{constante},
    \qquad
    \therefore \quad v_a=\text{constante}.
    \label{eq:fcbk-rapidez-areolar-constante}
\end{equation}
En un movimiento central, la rapidez areolar es constante.

Al parecer obtuvimos una nueva constante del movimiento, además del momentum angular. Sin embargo, la constancia de la rapidez areolar es una consecuencia de la constancia del momentum angular. En efecto,
\begin{equation}
    v_a=\frac{1}{2}r^2\dot\phi=\frac{1}{2}\frac{l}{m}=\frac{l}{2m},
    \label{eq:fcbk-va-l-sobre-2m}
\end{equation}
y por lo tanto
\begin{equation}
    \frac{dv_a}{dt}=\frac{1}{2m}\frac{dl}{dt}=0.
    \label{eq:fcbk-va-constante-por-l}
\end{equation}

De \eqref{eq:fcbk-derivada-r2phidot} y \eqref{eq:fcbk-rapidez-areolar} se tiene que
\begin{equation}
    a_\phi=2\dot r\dot\phi+r\ddot\phi
    =\frac{1}{r}\frac{d}{dt}(r^2\dot\phi)
    =\frac{1}{r}\frac{d}{dt}(2v_a)
    =\frac{2}{r}\frac{dv_a}{dt}=0.
    \label{eq:fcbk-aceleracion-angular-nula}
\end{equation}
El vector aceleración posee sólo componente radial:
\begin{equation}
    \vec a=(\ddot r-r\dot\phi^2)\hat r.
    \label{eq:fcbk-aceleracion-solo-radial}
\end{equation}

\subsection{Ecuaciones de itinerario}

De esta forma, se tiene que las ecuaciones del movimiento se simplifican a la forma
\begin{align}
    m(\ddot r-r\dot\phi^2)&=f(r),
    \label{eq:fcbk-edo-r-inicial}\\
    mr^2\dot\phi&=l.
    \label{eq:fcbk-edo-phi-inicial}
\end{align}
Esto es, una ecuación diferencial ordinaria de segundo orden en el tiempo para la variable $r$ y una ecuación diferencial ordinaria de primer orden en el tiempo para la variable $\phi$.

Se debe aclarar que las ecuaciones del movimiento originales son dos ecuaciones diferenciales ordinarias de segundo orden para $r$ y $\phi$; por lo tanto, al ser integradas debieran arrojar cuatro constantes de integración. Sin embargo, la conservación del momentum angular reduce el orden de la ecuación diferencial en $\phi$, lo que corresponde a una primera integral; por lo tanto, $l$ es una constante de integración, y sólo quedan tres constantes de integración.

La ecuación diferencial para $r$ se puede desacoplar de la variable $\phi$ y dejarla sólo en función de $r$. En efecto, de \eqref{eq:fcbk-edo-phi-inicial} se usa
\begin{equation}
    \dot\phi=\frac{l}{mr^2}.
    \label{eq:fcbk-phidot-l-mr2}
\end{equation}
Entonces, sustituyendo en \eqref{eq:fcbk-edo-r-inicial},
\begin{equation}
    \begin{aligned}
    m(\ddot r-r\dot\phi^2)&=f(r),\\
    m\left(\ddot r-r\left(\frac{l}{mr^2}\right)^2\right)&=f(r),\\
    m\left(\ddot r-\frac{l^2}{m^2r^3}\right)&=f(r).
    \end{aligned}
    \label{eq:fcbk-desacople-r}
\end{equation}
Por lo tanto,
\begin{equation}
    m\ddot r-\frac{l^2}{mr^3}=f(r).
    \label{eq:fcbk-ecuacion-r-desacoplada}
\end{equation}
Así, obtenemos las ecuaciones de itinerario del movimiento:
\begin{align}
    m\ddot r-\frac{l^2}{mr^3}&=f(r),
    \label{eq:fcbk-itinerario-r}\\
    mr^2\dot\phi&=l.
    \label{eq:fcbk-itinerario-phi}
\end{align}
Conocido $f(r)$ se puede obtener $r(t)$ de la primera ecuación y posteriormente de la segunda ecuación se obtiene $\phi(t)$.

\subsection{Ecuación de la trayectoria}

Sin embargo, las dos ecuaciones anteriores se pueden fundir en una sola. En efecto, de las ecuaciones del movimiento podemos conocer $r(t)$ y $\phi(t)$, pero podemos eliminar el tiempo de las ecuaciones despejando el tiempo en función del ángulo, $t=t(\phi)$, tal que
\begin{equation}
    r(t(\phi))=r(\phi),
    \label{eq:fcbk-r-de-phi}
\end{equation}
lo que se llama \textbf{ecuación de la trayectoria}.

Entonces,
\begin{equation}
    \begin{aligned}
    \dot r
    &=\frac{dr}{dt}\\
    &=\frac{d\phi}{dt}\frac{dr}{d\phi}\\
    &=\dot\phi\frac{dr}{d\phi}\\
    &=\frac{l}{mr^2}\frac{dr}{d\phi}\\
    &=-\frac{l}{m}\frac{d}{d\phi}\left(\frac{1}{r}\right).
    \end{aligned}
    \label{eq:fcbk-rdot-u}
\end{equation}
Luego,
\begin{equation}
    \begin{aligned}
    \ddot r
    &=\frac{d\dot r}{dt}\\
    &=\frac{d\phi}{dt}\frac{d\dot r}{d\phi}\\
    &=-\frac{l}{mr^2}\frac{l}{m}\frac{d^2}{d\phi^2}\left(\frac{1}{r}\right).
    \end{aligned}
    \label{eq:fcbk-rddot-u-intermedio}
\end{equation}
Por lo tanto,
\begin{equation}
    \ddot r=-\frac{l^2}{m^2r^2}\frac{d^2}{d\phi^2}\left(\frac{1}{r}\right).
    \label{eq:fcbk-rddot-u}
\end{equation}

Sustituyendo \eqref{eq:fcbk-rddot-u} en \eqref{eq:fcbk-itinerario-r}, se obtiene
\begin{equation}
    -\frac{l^2}{mr^2}\left[\frac{d^2}{d\phi^2}\left(\frac{1}{r}\right)+\frac{1}{r}\right]=f(r).
    \label{eq:fcbk-binet}
\end{equation}
Esta ecuación se llama \textbf{ecuación de Binet}. Logramos desacoplar las variables $r$ y $\phi$; sólo basta conocer $f(r)$.

Resumiendo las ecuaciones del movimiento:
\begin{align}
    -\frac{l^2}{mr^2}\left[\frac{d^2}{d\phi^2}\left(\frac{1}{r}\right)+\frac{1}{r}\right]&=f(r),
    \label{eq:fcbk-binet-resumen}\\
    \dot\phi&=\frac{l}{mr^2}.
    \label{eq:fcbk-phidot-resumen}
\end{align}
Conocido $f(r)$ podemos obtener $r(\phi)$, integrando la ecuación de Binet. Luego se introduce en la segunda ecuación y se puede obtener $\phi(t)$ y luego $r(t)$. No olvidar que hay tres constantes de integración.

\subsection{Potencial asociado a una fuerza central conservativa}

Podemos simplificar aún más la primera ecuación recordando que se impuso que la fuerza central sea conservativa, o sea existe un potencial asociado a la fuerza central:
\begin{equation}
    \vec F=-\vec\nabla U.
    \label{eq:fcbk-fuerza-potencial}
\end{equation}
Entonces,
\begin{equation}
    f(r)\hat r=-\frac{\partial U}{\partial r}\hat r-\frac{1}{r}\frac{\partial U}{\partial \phi}\hat\phi.
    \label{eq:fcbk-potencial-cilindrico}
\end{equation}
Claramente se observa que
\begin{equation}
    \frac{\partial U}{\partial \phi}=0,
    \label{eq:fcbk-U-no-depende-phi}
\end{equation}
y por lo tanto
\begin{equation}
    U=U(r).
    \label{eq:fcbk-U-de-r}
\end{equation}
Luego,
\begin{equation}
    \frac{dU}{dr}=-f(r),
    \label{eq:fcbk-dUdr}
\end{equation}
y entonces
\begin{equation}
    U(r)=-\int_{r_0}^{r}f(r)\,dr,
    \label{eq:fcbk-U-integral}
\end{equation}
donde $r_0$ es el radio de referencia donde se define $U(r_0)=0$.

Como la fuerza central es conservativa, entonces la energía mecánica se conserva:
\begin{equation}
    E=\frac{1}{2}mv^2+U(r).
    \label{eq:fcbk-energia-mecanica}
\end{equation}
La velocidad en coordenadas cilíndricas está dada por
\begin{equation}
    \vec v=\dot r\vec e_r+\dot\phi\vec e_\phi+\dot z\vec e_z=\dot r\hat r+r\dot\phi\hat\phi,
    \label{eq:fcbk-velocidad-cilindrica}
\end{equation}
y luego
\begin{equation}
    v^2=\vec v\cdot\vec v=\dot r^2+r^2\dot\phi^2.
    \label{eq:fcbk-v2-cilindrica}
\end{equation}
Por lo tanto, la energía se reescribe como
\begin{equation}
    E=\frac{1}{2}m(\dot r^2+r^2\dot\phi^2)+U(r).
    \label{eq:fcbk-energia-r-phi}
\end{equation}
Usando \eqref{eq:fcbk-phidot-l-mr2}, se obtiene
\begin{equation}
    E=\frac{1}{2}m\dot r^2+\frac{l^2}{2mr^2}+U(r).
    \label{eq:fcbk-energia-reducida}
\end{equation}
Tenemos dos ecuaciones de primer orden:
\begin{align}
    E&=\frac{1}{2}m\dot r^2+\frac{l^2}{2mr^2}+U(r),
    \label{eq:fcbk-energia-reducida-resumen}\\
    l&=mr^2\dot\phi.
    \label{eq:fcbk-l-resumen-final}
\end{align}
Se ha reducido el orden de la derivada temporal en la ecuación radial en un orden, por lo que la energía es una constante de integración.

\newpage
\section{Problema de Kepler}

\subsection{Ley de gravitación universal y ecuación de Binet}

La idea es estudiar el movimiento de una partícula sometida a la acción de una fuerza central dependiente del inverso del cuadrado del radio, esto es, la \textbf{ley de gravitación universal de Newton}:
\begin{equation}
    \vec F=-\frac{k}{r^2}\hat r,
    \qquad
    f(r)=-\frac{k}{r^2}.
    \label{eq:fcbk-kepler-fuerza-inversa-cuadrado}
\end{equation}
Introduciéndola en la ecuación de Binet \eqref{eq:fcbk-binet}, se tiene
\begin{equation}
    -\frac{l^2}{mr^2}\left[\frac{d^2}{d\phi^2}\left(\frac{1}{r}\right)+\frac{1}{r}\right]
    =-\frac{k}{r^2}.
    \label{eq:fcbk-kepler-binet-sustitucion}
\end{equation}
Por lo tanto,
\begin{equation}
    \frac{d^2}{d\phi^2}\left(\frac{1}{r}\right)+\frac{1}{r}=\frac{km}{l^2}.
    \label{eq:fcbk-kepler-binet-final}
\end{equation}

Si efectuamos el cambio de variable
\begin{equation}
    u=\frac{1}{r},
    \label{eq:fcbk-kepler-u}
\end{equation}
entonces la ecuación de Binet se simplifica en la forma
\begin{equation}
    \frac{d^2u}{d\phi^2}+u=\frac{km}{l^2}.
    \label{eq:fcbk-kepler-u-edo}
\end{equation}
La solución de esta ecuación es dada por
\begin{equation}
    u=u_p+u_h,
    \label{eq:fcbk-kepler-u-general}
\end{equation}
donde $u_p$ es la solución particular,
\begin{equation}
    u_p=\frac{km}{l^2},
    \label{eq:fcbk-kepler-up}
\end{equation}
y $u_h$ corresponde a la solución de la ecuación homogénea
\begin{equation}
    \frac{d^2u_h}{d\phi^2}+u_h=0,
    \label{eq:fcbk-kepler-uh-edo}
\end{equation}
cuya solución es
\begin{equation}
    u_h=A\sin(\phi-\phi_0),
    \label{eq:fcbk-kepler-uh}
\end{equation}
donde $A$ y $\phi_0$ son constantes de integración. La solución a la ecuación de Binet entonces es
\begin{equation}
    u(\phi)=A\sin(\phi-\phi_0)+\frac{km}{l^2}.
    \label{eq:fcbk-kepler-u-sol}
\end{equation}
Equivalentemente,
\begin{equation}
    \frac{1}{r}=\frac{km}{l^2}\left(1+\frac{Al^2}{mk}\sin(\phi-\phi_0)\right).
    \label{eq:fcbk-kepler-inversa-r}
\end{equation}
Así,
\begin{equation}
    r(\phi)=\frac{\frac{l^2}{km}}{1+\frac{Al^2}{mk}\sin(\phi-\phi_0)}.
    \label{eq:fcbk-kepler-r-sol-intermedia}
\end{equation}
Si definimos
\begin{equation}
    p=\frac{l^2}{km},
    \qquad
    e=\frac{Al^2}{mk},
    \label{eq:fcbk-kepler-p-e-def}
\end{equation}
tenemos
\begin{equation}
    r(\phi)=\frac{p}{1+e\sin(\phi-\phi_0)}.
    \label{eq:fcbk-kepler-r-conica-general}
\end{equation}
Si ajustamos de forma apropiada el eje polar tenemos $\phi_0=0$; por lo tanto,
\begin{equation}
    r(\phi)=\frac{p}{1+e\sin\phi},
    \label{eq:fcbk-kepler-r-conica}
\end{equation}
la que corresponde a la ecuación focal de una sección cónica de parámetro $p$ y excentricidad $e$.

\subsection{Secciones cónicas}

Las soluciones del problema de Kepler se organizan geométricamente en términos de \textbf{secciones cónicas}: circunferencia, elipse, parábola e hipérbola. Estas formas se muestran en la Figura~\ref{fig:fcbk-secciones-conicas}.

\begin{figure}[H]
    \centering
    \IfFileExists{figures/09_fuerzas_centrales_binet_kepler/secciones_conicas.png}{%
        \includegraphics[width=0.7\linewidth]{figures/09_fuerzas_centrales_binet_kepler/secciones_conicas.png}%
    }{%
        \fbox{\parbox{0.7\linewidth}{\centering Figura pendiente: secciones cónicas.}}%
    }
    \caption{Secciones cónicas asociadas a distintos valores de la excentricidad.}
    \label{fig:fcbk-secciones-conicas}
\end{figure}

La ecuación anterior posee una constante de integración $A$, la que podemos ajustar con la energía del sistema. En efecto, la energía potencial del sistema está dada por
\begin{equation}
    U(r)=-\int_{\infty}^{r}\frac{k}{r^2}\,dr=-\frac{k}{r},
    \label{eq:fcbk-kepler-potencial}
\end{equation}
donde el radio de referencia se fijó en infinito, puesto que la fuerza tiende a cero en infinito. Entonces
\begin{equation}
    E=\frac{1}{2}m\dot r^2+\frac{l^2}{2mr^2}-\frac{k}{r}.
    \label{eq:fcbk-kepler-energia}
\end{equation}

Despejando $\dot r$, se obtiene
\begin{equation}
    \dot r^2=\frac{2}{m}\left(E-\frac{l^2}{2mr^2}+\frac{k}{r}\right).
    \label{eq:fcbk-kepler-rdot2}
\end{equation}
Pero, por \eqref{eq:fcbk-rdot-u},
\begin{equation}
    \dot r=\frac{l}{mr^2}\frac{dr}{d\phi}.
    \label{eq:fcbk-kepler-rdot-dr-dphi}
\end{equation}
Sea $u=1/r$. Al integrar se tiene nuevamente
\begin{equation}
    r(\phi)=\frac{p}{1+e\sin(\phi-\phi_0)},
    \label{eq:fcbk-kepler-r-integrada}
\end{equation}
con
\begin{equation}
    p=\frac{l^2}{km},
    \qquad
    e=\sqrt{1+\frac{2El^2}{mk^2}}.
    \label{eq:fcbk-kepler-excentricidad-energia}
\end{equation}
La expresión \eqref{eq:fcbk-kepler-r-integrada} corresponde a la ecuación focal de una sección cónica de parámetro $p$ y excentricidad $e$.

\subsection{Excentricidad como razón foco-directriz}

Sean $F$ un punto fijo, llamado \textbf{foco}, y una recta fija, llamada \textbf{directriz}, en el plano. Sea $P$ un punto sobre la curva en el plano y se define la excentricidad $e$ como el cociente entre la distancia del punto $P$ al foco y la distancia del punto $P$ a la directriz:
\begin{equation}
    e=\frac{PF}{PQ}=\frac{r}{PQ}.
    \label{eq:fcbk-excentricidad-foco-directriz}
\end{equation}
La geometría foco-directriz de una cónica se muestra en la Figura~\ref{fig:fcbk-conica-foco-directriz}.

\begin{figure}[H]
    \centering
    \IfFileExists{figures/09_fuerzas_centrales_binet_kepler/conica_foco_directriz.png}{%
        \includegraphics[width=0.45\linewidth]{figures/09_fuerzas_centrales_binet_kepler/conica_foco_directriz.png}%
    }{%
        \fbox{\parbox{0.45\linewidth}{\centering Figura pendiente: foco, directriz y excentricidad.}}%
    }
    \caption{Definición geométrica de la excentricidad mediante foco y directriz.}
    \label{fig:fcbk-conica-foco-directriz}
\end{figure}

Sea $d$ la distancia entre el foco, en el polo, y la directriz correspondiente. Entonces,
\begin{equation}
    r\sin\phi+PQ=d.
    \label{eq:fcbk-conica-relacion-d}
\end{equation}
Usando \eqref{eq:fcbk-excentricidad-foco-directriz}, se tiene $PQ=r/e$, de modo que
\begin{equation}
    r\sin\phi+\frac{r}{e}=d.
    \label{eq:fcbk-conica-relacion-pq}
\end{equation}
Por lo tanto,
\begin{equation}
    r(e\sin\phi+1)=ed,
    \label{eq:fcbk-conica-intermedia}
\end{equation}
y finalmente
\begin{equation}
    r(\phi)=\frac{ed}{1+e\sin\phi}.
    \label{eq:fcbk-conica-focal}
\end{equation}
Si la distancia $d$ es constante, entonces la ecuación anterior es la ecuación focal, medida desde el foco, de una sección cónica. La excentricidad $e$ clasifica el tipo de cónica:
\begin{itemize}
    \item la cónica es una \textbf{circunferencia} si $e=0$;
    \item la cónica es una \textbf{elipse} si $0<e<1$;
    \item la cónica es una \textbf{parábola} si $e=1$;
    \item la cónica es una \textbf{hipérbola} si $e>1$.
\end{itemize}

\subsection{Energía y clasificación de órbitas}

Volviendo a la energía, es claro que la energía está acotada. Desde \eqref{eq:fcbk-kepler-excentricidad-energia},
\begin{equation}
    e=\sqrt{1+\frac{2El^2}{mk^2}}.
    \label{eq:fcbk-e-raiz-energia}
\end{equation}
Como el argumento de la raíz debe ser no negativo,
\begin{equation}
    1+\frac{2El^2}{mk^2}\geq 0,
    \label{eq:fcbk-condicion-raiz-e}
\end{equation}
se concluye que
\begin{equation}
    E\geq -\frac{mk^2}{2l^2}.
    \label{eq:fcbk-cota-energia}
\end{equation}
Por lo tanto, existe una energía mínima
\begin{equation}
    E_{\min}=-\frac{mk^2}{2l^2},
    \label{eq:fcbk-energia-minima}
\end{equation}
la que es negativa.

De acuerdo a los valores de $e$, tenemos que la trayectoria puede ser hiperbólica, parabólica, elíptica o circular:
\begin{itemize}
    \item si $e>1$, entonces $E>0$ y la trayectoria es una \textbf{hipérbola};
    \item si $e=1$, entonces $E=0$ y la trayectoria es una \textbf{parábola};
    \item si $0<e<1$, entonces $E_{\min}<E<0$ y la trayectoria es una \textbf{elipse};
    \item si $e=0$, entonces $E=E_{\min}$ y la trayectoria es una \textbf{circunferencia} de radio
    \begin{equation}
        r=p=\frac{l^2}{km}=\text{constante}.
        \label{eq:fcbk-radio-circular}
    \end{equation}
\end{itemize}

En resumen:
\begin{center}
\begin{tabular}{c c c}
\hline
$e$ & Órbita & Energía \\
\hline
$e>1$ & Hipérbola & $E>0$ \\
$e=1$ & Parábola & $E=0$ \\
$0<e<1$ & Elipse & $E_{\min}<E<0$ \\
$e=0$ & Circunferencia & $E=E_{\min}=-\dfrac{mk^2}{2l^2}$ \\
\hline
\end{tabular}
\end{center}

Si $E\geq 0$, entonces $T\geq |U|$ y aparecen hipérbolas y parábolas. Si $E<0$, entonces $T\leq |U|$ y aparecen elipses y circunferencias.

\subsection{Leyes de Kepler}

Las \textbf{leyes de Kepler} pueden escribirse como:
\begin{itemize}
    \item \textbf{Primera ley:} los planetas describen órbitas elípticas en torno del Sol, y el Sol se encuentra en uno de los focos de la elipse.
    \item \textbf{Segunda ley:} el radio vector barre áreas iguales en tiempos iguales; la rapidez areolar es constante.
    \item \textbf{Tercera ley:} el cuadrado del periodo de la órbita es proporcional al cubo del semieje mayor:
    \begin{equation}
        T^2=\frac{4\pi^2}{GM}a^3.
        \label{eq:fcbk-tercera-ley-kepler}
    \end{equation}
\end{itemize}

La Tierra gira alrededor de una estrella llamada Sol, en una órbita elíptica con parámetros: excentricidad $0{,}01671123$, perihelio $147098290\,\mathrm{km}$, afelio $152098232\,\mathrm{km}$ e inclinación $23{,}5^\circ$. La Figura~\ref{fig:fcbk-orbita-tierra} muestra una representación de la órbita terrestre.

\begin{figure}[H]
    \centering
    \IfFileExists{figures/09_fuerzas_centrales_binet_kepler/orbita_tierra_parametros.png}{%
        \includegraphics[width=0.62\linewidth]{figures/09_fuerzas_centrales_binet_kepler/orbita_tierra_parametros.png}%
    }{%
        \fbox{\parbox{0.62\linewidth}{\centering Figura pendiente: órbita terrestre.}}%
    }
    \caption{Representación de la órbita terrestre con perihelio, afelio, equinoccios y solsticios.}
    \label{fig:fcbk-orbita-tierra}
\end{figure}

Es conveniente insistir que en un movimiento central el momentum angular con respecto a cualquier eje que pasa por el centro de fuerzas es constante. La importancia del momentum angular se debe precisamente al gran número de fuerzas centrales que aparecen en la física.

Es conveniente usar coordenadas polares $(r,\phi)$ y la fórmula de Binet en los casos en que $r$ y $\phi$ están relacionados de una forma sencilla:
\begin{equation}
    -\frac{l^2}{mr^2}\left[\frac{d^2}{d\phi^2}\left(\frac{1}{r}\right)+\frac{1}{r}\right]=f(r).
    \label{eq:fcbk-binet-uso-practico}
\end{equation}
Cuando $r(\phi)$ es una función complicada, es conveniente usar coordenadas polares $(r,\phi)$ y el principio de conservación de energía:
\begin{equation}
    \frac{l^2}{2mr^4}\left[\left(\frac{dr}{d\phi}\right)^2+r^2\right]+U(r)=E.
    \label{eq:fcbk-energia-r-de-phi}
\end{equation}
Si se conocen las condiciones iniciales $l$ y $E$, se puede, en principio, determinar $r=r(\phi)$.

\newpage
\section{Problema unidimensional equivalente}

\subsection{Potencial efectivo}

Volvamos a la ecuación de energía
\begin{equation}
    E=\frac{1}{2}m\dot r^2+\frac{l^2}{2mr^2}+U(r).
    \label{eq:fcbk-energia-para-potencial-efectivo}
\end{equation}
A partir de esta ecuación podemos obtener información cualitativa importante.

Si definimos el nuevo potencial
\begin{equation}
    V_{\mathrm{ef}}(r)=\frac{l^2}{2mr^2}+U(r),
    \label{eq:fcbk-potencial-efectivo}
\end{equation}
llamado \textbf{potencial efectivo}, entonces la energía se escribe
\begin{equation}
    E=\frac{1}{2}m\dot r^2+V_{\mathrm{ef}}(r).
    \label{eq:fcbk-energia-potencial-efectivo}
\end{equation}

El término
\begin{equation}
    V_c(r)=\frac{l^2}{2mr^2}
    \label{eq:fcbk-potencial-centrifugo}
\end{equation}
en $V_{\mathrm{ef}}(r)$ es la energía potencial asociada a la fuerza centrífuga y es llamada \textbf{potencial centrífugo}. En efecto,
\begin{equation}
    F_c=-\frac{dV_c}{dr}
    =-\frac{d}{dr}\left(\frac{l^2}{2mr^2}\right)
    =\frac{l^2}{mr^3}.
    \label{eq:fcbk-fuerza-centrifuga-potencial}
\end{equation}
Usando $l=mr^2\dot\phi$, se tiene
\begin{equation}
    F_c=\frac{(mr^2\dot\phi)^2}{mr^3}=mr\dot\phi^2=mr\omega^2=m\frac{v^2}{r}.
    \label{eq:fcbk-fuerza-centrifuga-formas}
\end{equation}
Esta fuerza no apunta hacia el centro, sino hacia afuera. Esta es la razón del nombre de $V_c$.

\subsection{Región de movimiento y puntos de retorno}

Los límites de la región del movimiento son los valores de $r$ para los cuales se cumple
\begin{equation}
    \frac{1}{2}m\dot r^2=E-V_{\mathrm{ef}}(r)\geq 0.
    \label{eq:fcbk-region-movimiento}
\end{equation}
En el límite de la región, $E=V_{\mathrm{ef}}(r)$, se cumple $\dot r=0$, lo que define un \textbf{punto de retorno} de la trayectoria.

Si existe un valor mínimo de $r$, es decir, si existe un $r_{\min}$ tal que $r\geq r_{\min}$, entonces el movimiento es infinito, o sea, el movimiento comienza y termina en el infinito; la Figura~\ref{fig:fcbk-radio-minimo-orbita-abierta} representa esta situación.

\begin{figure}[H]
    \centering
    \IfFileExists{figures/09_fuerzas_centrales_binet_kepler/radio_minimo_orbita_abierta.png}{%
        \includegraphics[width=0.50\linewidth]{figures/09_fuerzas_centrales_binet_kepler/radio_minimo_orbita_abierta.png}%
    }{%
        \fbox{\parbox{0.50\linewidth}{\centering Figura pendiente: radio mínimo.}}%
    }
    \caption{Trayectoria no acotada con un radio mínimo $r_{\min}$.}
    \label{fig:fcbk-radio-minimo-orbita-abierta}
\end{figure}

Si existen $r_{\min}$ y $r_{\max}$, es decir, $r_{\min}\leq r\leq r_{\max}$, entonces el movimiento es finito, lo que significa que el movimiento es acotado. El análisis cualitativo se lleva a cabo sobre un gráfico de $V_{\mathrm{ef}}$ versus $r$, que se construye trazando separadamente las curvas $-k/r$ y $l^2/(2mr^2)$, como en la Figura~\ref{fig:fcbk-potencial-efectivo-componentes}.

\begin{figure}[H]
    \centering
    \IfFileExists{figures/09_fuerzas_centrales_binet_kepler/potencial_efectivo_componentes.png}{%
        \includegraphics[width=0.52\linewidth]{figures/09_fuerzas_centrales_binet_kepler/potencial_efectivo_componentes.png}%
    }{%
        \fbox{\parbox{0.52\linewidth}{\centering Figura pendiente: componentes del potencial efectivo.}}%
    }
    \caption{Construcción del potencial efectivo a partir de $-k/r$ y $l^2/(2mr^2)$.}
    \label{fig:fcbk-potencial-efectivo-componentes}
\end{figure}

Si sumamos estas dos curvas obtenemos $V_{\mathrm{ef}}$. El potencial efectivo posee un cero en
\begin{equation}
    r_c=\frac{l^2}{2km},
    \label{eq:fcbk-rc-potencial-efectivo}
\end{equation}
y además posee un mínimo ubicado en
\begin{equation}
    r_{\min}=\frac{l^2}{km}=2r_c,
    \label{eq:fcbk-rmin-potencial-efectivo}
\end{equation}
tal que
\begin{equation}
    V_{\mathrm{ef,min}}=-\frac{k^2m}{2l^2}=E_{\min}.
    \label{eq:fcbk-vef-min}
\end{equation}
La forma cualitativa de $V_{\mathrm{ef}}$ se muestra en la Figura~\ref{fig:fcbk-potencial-efectivo-minimo}.

\begin{figure}[H]
    \centering
    \IfFileExists{figures/09_fuerzas_centrales_binet_kepler/potencial_efectivo_minimo.png}{%
        \includegraphics[width=0.55\linewidth]{figures/09_fuerzas_centrales_binet_kepler/potencial_efectivo_minimo.png}%
    }{%
        \fbox{\parbox{0.55\linewidth}{\centering Figura pendiente: mínimo del potencial efectivo.}}%
    }
    \caption{Potencial efectivo con cero en $r_c$ y mínimo en $r_{\min}$.}
    \label{fig:fcbk-potencial-efectivo-minimo}
\end{figure}

A partir de este gráfico es posible obtener una descripción cualitativa de las posibles trayectorias. La barrera centrífuga $V_c(r)$ impide que las partículas caigan al centro de fuerzas cualquiera que sea la energía, siempre y cuando $l\neq 0$.

\subsection{Regiones permitidas y prohibidas}

A partir del gráfico $V_{\mathrm{ef}}$ versus $r$ podemos trazar líneas de energía constante. El movimiento no es posible si
\begin{equation}
    \frac{1}{2}m\dot r^2<0,
    \label{eq:fcbk-region-prohibida-condicion}
\end{equation}
por lo que se define una \textbf{región prohibida} para el movimiento de una partícula, como se indica en la Figura~\ref{fig:fcbk-region-prohibida}.

\begin{figure}[H]
    \centering
    \IfFileExists{figures/09_fuerzas_centrales_binet_kepler/energia_constante_region_prohibida.png}{%
        \includegraphics[width=0.55\linewidth]{figures/09_fuerzas_centrales_binet_kepler/energia_constante_region_prohibida.png}%
    }{%
        \fbox{\parbox{0.55\linewidth}{\centering Figura pendiente: región prohibida.}}%
    }
    \caption{Línea de energía constante y región prohibida para el movimiento radial.}
    \label{fig:fcbk-region-prohibida}
\end{figure}

Consideremos una partícula de energía $E_1>0$. Entonces $E_1-V_{\mathrm{ef}}(r)\geq 0$ y en el punto de retorno $r=r_A$ se tiene $E_1=V_{\mathrm{ef}}(r_A)$. Cuando $E_1\geq V_{\mathrm{ef}}(r)$ se tiene $r_A\leq r<\infty$. Tenemos una región permitida y una región prohibida, como se muestra en la Figura~\ref{fig:fcbk-energia-positiva-retorno}.

\begin{figure}[H]
    \centering
    \IfFileExists{figures/09_fuerzas_centrales_binet_kepler/energia_positiva_punto_retorno.png}{%
        \includegraphics[width=0.55\linewidth]{figures/09_fuerzas_centrales_binet_kepler/energia_positiva_punto_retorno.png}%
    }{%
        \fbox{\parbox{0.55\linewidth}{\centering Figura pendiente: energía positiva y punto de retorno.}}%
    }
    \caption{Caso $E_1>0$: existe un punto de retorno inferior y el movimiento no está acotado superiormente.}
    \label{fig:fcbk-energia-positiva-retorno}
\end{figure}

La analogía con regiones permitidas y prohibidas se ilustra en la Figura~\ref{fig:fcbk-regiones-permitidas-prohibidas}: la partícula no puede entrar en regiones donde la energía cinética radial sería negativa.

\begin{figure}[H]
    \centering
    \IfFileExists{figures/09_fuerzas_centrales_binet_kepler/regiones_permitidas_prohibidas.png}{%
        \includegraphics[width=0.65\linewidth]{figures/09_fuerzas_centrales_binet_kepler/regiones_permitidas_prohibidas.png}%
    }{%
        \fbox{\parbox{0.65\linewidth}{\centering Figura pendiente: regiones permitidas y prohibidas.}}%
    }
    \caption{Analogía cualitativa entre regiones permitidas y prohibidas del movimiento.}
    \label{fig:fcbk-regiones-permitidas-prohibidas}
\end{figure}

De la figura es claro que la partícula nunca podrá estar a una distancia menor que $r_A$ del origen, ya que de otro modo $m\dot r^2/2<0$, por lo que $\dot r$ correspondería a una rapidez radial imaginaria. Por otro lado, el valor de $r$ no tiene límite superior, por lo cual la órbita no está acotada y por lo tanto no está cerrada. Una partícula procedente del infinito rebota en la barrera centrífuga dirigiéndose de nuevo hacia el infinito.

Consideremos ahora una partícula de energía $E_2=0$. En este caso $m\dot r^2/2=V_{\mathrm{ef}}$ y en el punto de retorno $r=r_c$ se tiene $\dot r=0$. Por lo tanto, el movimiento es infinito, $r_c\leq r<\infty$, como se muestra en la Figura~\ref{fig:fcbk-energia-cero-retorno}.

\begin{figure}[H]
    \centering
    \IfFileExists{figures/09_fuerzas_centrales_binet_kepler/energia_cero_punto_retorno.png}{%
        \includegraphics[width=0.55\linewidth]{figures/09_fuerzas_centrales_binet_kepler/energia_cero_punto_retorno.png}%
    }{%
        \fbox{\parbox{0.55\linewidth}{\centering Figura pendiente: energía cero y punto de retorno.}}%
    }
    \caption{Caso $E_2=0$: el movimiento es infinito y posee un punto de retorno.}
    \label{fig:fcbk-energia-cero-retorno}
\end{figure}

Consideremos una partícula de energía $E_3<0$. En este caso se obtiene que además del límite inferior en el radio, $r_B$, también existe un valor máximo para el radio, $r_D$, que no puede sobrepasar con energía cinética radial positiva. El movimiento está limitado y existen dos puntos de retorno, como se observa en la Figura~\ref{fig:fcbk-energia-negativa-retorno}.

\begin{figure}[H]
    \centering
    \IfFileExists{figures/09_fuerzas_centrales_binet_kepler/energia_negativa_puntos_retorno.png}{%
        \includegraphics[width=0.55\linewidth]{figures/09_fuerzas_centrales_binet_kepler/energia_negativa_puntos_retorno.png}%
    }{%
        \fbox{\parbox{0.55\linewidth}{\centering Figura pendiente: energía negativa y puntos de retorno.}}%
    }
    \caption{Caso $E_3<0$: el movimiento queda acotado entre dos puntos de retorno.}
    \label{fig:fcbk-energia-negativa-retorno}
\end{figure}

Esto no significa que las órbitas tengan que ser necesariamente cerradas. Lo más que cabe asegurar es que las órbitas están limitadas por dos circunferencias de radios $r_B$ y $r_D$, estando los puntos absidales en dichas circunferencias, como se ilustra en la Figura~\ref{fig:fcbk-orbita-acotada-absidales}.

\begin{figure}[H]
    \centering
    \IfFileExists{figures/09_fuerzas_centrales_binet_kepler/orbita_acotada_absidales.png}{%
        \includegraphics[width=0.52\linewidth]{figures/09_fuerzas_centrales_binet_kepler/orbita_acotada_absidales.png}%
    }{%
        \fbox{\parbox{0.52\linewidth}{\centering Figura pendiente: órbita acotada y puntos absidales.}}%
    }
    \caption{Órbita acotada entre dos circunferencias de radios $r_B$ y $r_D$.}
    \label{fig:fcbk-orbita-acotada-absidales}
\end{figure}

Finalmente, consideremos una partícula de energía $E_4=E_{\min}$. Este es precisamente el mínimo de la curva de potencial efectivo: los límites mínimo y máximo coinciden. En este caso sólo es posible que
\begin{equation}
    \dot r=0,
    \qquad
    r=r_{\min}=\text{constante},
    \label{eq:fcbk-orbita-circular-condicion}
\end{equation}
por lo tanto, la órbita es una circunferencia. La Figura~\ref{fig:fcbk-orbita-circular-energia-minima} muestra este caso.

\begin{figure}[H]
    \centering
    \IfFileExists{figures/09_fuerzas_centrales_binet_kepler/orbita_circular_energia_minima.png}{%
        \includegraphics[width=0.55\linewidth]{figures/09_fuerzas_centrales_binet_kepler/orbita_circular_energia_minima.png}%
    }{%
        \fbox{\parbox{0.55\linewidth}{\centering Figura pendiente: órbita circular.}}%
    }
    \caption{Caso $E_4=E_{\min}$: el movimiento circular ocurre en el mínimo del potencial efectivo.}
    \label{fig:fcbk-orbita-circular-energia-minima}
\end{figure}

Si la energía satisface
\begin{equation}
    E<E_{\min},
    \label{eq:fcbk-energia-menor-minima}
\end{equation}
el movimiento no es posible, puesto que la velocidad radial sería imaginaria, $\dot r\in\mathbb{C}$.

\end{document}
